% status: ZERO
% Chapter 14 of the second edition. Plan: docs/STRUCTURE.md. Rules: AUTHORING.md.
% The sections from "One prediction is not yet a service" to "What we have
% established" are the first edition's Chapter 1, moved here with git mv on
% 2026-09-23. They are the mechanism of this chapter; the ZERO pass reworks
% them around the anatomy. Commands are regression checked by
% scripts/check-textbook.py.
% Measurements: research/measurements/2026-09-24-m1-part3.md (lifecycle, cancel2,
% sse_bytes).
\chapter{The Request Lifecycle: Streaming, Cancellation and Terminal Ownership}\label{ch:request-lifecycle}\label{chapter-1-the-missing-half-of-ai}

\begin{ChapterIntent}
A request is not a function call. It is a state machine that receives tokens
from a model, turns them into bytes a client can read, honours cancellation
and budgets, survives failures part-way through a stream, and ends with
exactly one terminal outcome. This chapter makes that contract explicit,
proves the single-owner property that keeps it true, and defines the
timestamps --- arrival, first token, first text, completion --- from which
time to first token and inter-token latency are measured. It builds on the
first edition's executable request trace.
\end{ChapterIntent}

\OpeningSection{Time to first token, measured}

Ask llama.cpp's server on the M1 for 128 tokens after a 200-token prompt, and
stream the answer. In five runs the first token arrived \textbf{741 to 762
milliseconds} after the request was sent; after that a token arrived every 40
milliseconds or so --- median gaps of 39.1 to 40.5~ms, with the slowest gap of
each run between 53 and 83~ms --- and the whole answer took about six seconds
(measured with Llama~3.2~3B; record \texttt{2026-09-24-m1-part3}). The server's
own accounting says where those first three-quarters of a second went: 738 to
757~ms of processing the prompt. Everything else --- parsing the request, scheduling it, sampling,
detokenizing, framing an event, crossing a socket --- fit in the few
milliseconds left.

Then the clients hung up. A client that closed its streaming connection after
its twentieth token left a server that noticed within 71~ms: the request's slot
was released after 21 tokens, one of them generated after the disconnect. A
client that asked for a long answer \emph{without} streaming and hung up after
three seconds left a server that behaved in two ways. On an otherwise idle
server, the request was cancelled within a second and a half. While another
client was streaming, it was not cancelled at all: the server kept generating
for the absent client for another 37 seconds, some 490 tokens nobody would
read, and stopped only when the other stream ended. The case study at the end
of this chapter explains why.

A request is not a function call. It is a small state machine with a clock
attached, and its contract --- what the client sees, when, what the engine
spends, and when it stops --- decides both the latency users feel and the
capacity the engine wastes. This chapter makes that contract explicit.

\NapkinSection{Napkin math: what a stream costs the engine}

A streamed token becomes an event on the wire. In the measured build each
event on llama.cpp's own completion endpoint was about 122 bytes and on its
OpenAI-compatible chat endpoint about 306 bytes, for a token that carries three
or four characters of text (measured). At 25 tokens per second a stream sends
about 7.7~KB/s on the OpenAI-compatible endpoint, and ten thousand concurrent
streams send about 77~MB/s: real, but small next to the machines producing
them (derived). Each event also costs CPU for detokenization and
serialization, which engines move off the scheduling loop
(\Chref{continuous-batching}).

The larger cost of a stream is its duration. From admission to its terminal
event a request holds a batch slot and its KV cache: for Qwen3-8B at 4,096
tokens of context, 576~MiB (\Chref{kv-cache}). A request that should have been
cancelled but runs to completion wastes both, and the waste is easy to price:
\begin{equation}
W_{\text{orphan}} = n_{\text{left}}\;t_{\text{token}}\quad\text{of a slot, holding }m_{kv}\,S\text{ bytes},
\label{eq:orphan-waste}
\end{equation}
for $n_{\text{left}}$ tokens left to generate at $t_{\text{token}}$ each. The
opening's orphan held a slot for 37 seconds. At a hosted service's scale, a
fraction of a percent of abandoned requests that are never cancelled is a
fraction of a percent of the fleet serving nobody.

\section{One prediction is not yet a service}
Suppose a model has just produced four numbers. One number is larger than the
others. What must happen before a user can read an answer?

Someone must interpret the index of that number, decide whether it means text
or termination, turn it into bytes, preserve their order, update the input for
the next prediction and release the request's resources. If the user cancels
halfway through, the computation must stop with an honest outcome. If the model
fails after emitting a prefix, that prefix does not become a successful answer
merely because some bytes reached the client.

This is the subject of an \Term{inference engine}: the running system around
the numerical prediction. Our first object will be one request, not an
architecture poster. We will run it, draw its state, change its stopping
condition and identify which component owns the final outcome.

\Callout{The boundary of this chapter}{The companion executable already contains
the four-token numerical fixture derived in \Appref{smallest-model}. Here we treat its
forward calculation as an opaque function and inspect the lifecycle around it.
We do not claim to have built a Transformer, trained a model or implemented a
network server. This keeps the first experiment reproducible on the current
checkout without relying on an obsolete fake-model API.}

\section{Run a request you can account for}
From the repository root, run:
\begin{verbatim}
cd code/mini-engine
cargo run -q -p engine0 -- --trace 'I like'
\end{verbatim}
The fixture vocabulary has four IDs: 0 is end of sequence, 1 is \texttt{I},
2 is \texttt{\textvisiblespace like}, and 3 is
\texttt{\textvisiblespace Rust}. The visible-space symbol makes the leading
byte explicit; it is not part of the spelling typed at the terminal.
These are identities, not quantities. The prompt becomes $[1,2]$.
The first forward pass uses its final ID, 2, and returns the scores
\begin{equation}
z^{(0)}=(-0.7,\;0.1,\;0.4,\;2.2).
\label{eq:first-logits}
\end{equation}
The largest score belongs to ID 3. Greedy selection therefore chooses
\texttt{Rust}. The next forward pass receives the prompt plus that new ID,
and its largest score belongs to ID 0. This time selection means stop.
The generated bytes spell \texttt{\textvisiblespace Rust}; the executable does not echo the
prompt as generated text.

Do not memorize the four numbers yet. \Appref{smallest-model} derives every multiplication.
For now, check three different observations in the trace: a model invocation,
a selected ID and a terminal outcome. They are not interchangeable.
There are two model invocations but only one emitted text token. The selected
EOS control ID ends the request; it is not appended as visible text.

The printed microsecond values will vary. They include instrumentation and
local runtime overhead. They are not benchmark results and should not be copied
into a performance comparison.

\EngineFigure{request-trace}{One request through two predictions. The solid
arrows carry history or a chosen ID; the dashed arrow is a stopping decision.
EOS is selected but is not a visible output token.}{fig:request-trace}

Now change only the prompt to \texttt{I}. Its first logits are approximately
$(0.2,0.2,0.2,0.1)$. Greedy selection breaks the leading tie by the smallest ID,
so the request ends immediately. No text event is required for a valid
completion. This small counterexample is useful: a demonstration that promises
a word for every prompt has already imposed a behavior the model does not have.

\section{Make the loop explicit}
Let $x$ be the input sequence of token IDs and let $g_t$ be the list of
\emph{visible generated tokens} before iteration $t$. The model $F_\theta$,
with fixed parameters $\theta$, produces a vector of logits:
\begin{equation}
z_t=F_\theta(x\mathbin{\|}g_t),\qquad
a_t=\operatorname*{arg\,max}_{i}z_{t,i}.
\label{eq:lifecycle-choice}
\end{equation}
The symbol $\|$ denotes concatenation. The vector index $i$ ranges over the
model vocabulary. The argmax needs a tie rule; our deterministic rule prefers
the smallest ID. Scores need not be positive, normalized or interpretable as
confidence.

For an ordinary text token, the next state is
\begin{equation}
g_{t+1}=g_t\mathbin{\|}(a_t).
\label{eq:lifecycle-feedback}
\end{equation}
For EOS, there is no next generation state: the loop selects a terminal
outcome. A maximum-token budget can terminate the loop before another model
invocation. Cancellation is checked at declared boundaries. These details
belong to the runtime policy rather than to the matrix calculation.

This recurrence is \Term{autoregression}. It says what feeds the next call;
it does not say that the model actually uses every earlier ID. Our fixture
looks only at the final input token. Later attention chapters will construct
genuine dependence on earlier positions. Keeping the history and using that
history are separate claims.

Here is the conceptual algorithm. It is pseudocode, not a second implementation
that silently replaces the tested Rust:
\begin{verbatim}
validate request and initialize request-local state
repeat:
    if cancellation is observed: choose Cancelled; leave loop
    if generated count reaches budget: choose MaxTokens; leave loop
    logits = model.forward(prompt + generated)
    if forward fails: choose Failed; leave loop
    validate logits and apply the request's selection policy
    if selected ID is EOS: choose EndOfSequence; leave loop
    append the selected text ID to request history
    send its token event
    decode its bytes and send any complete text prefix
send exactly one terminal outcome
\end{verbatim}
The implementation also checks cancellation after the model returns and
validates incomplete UTF-8 at completion. Error handling is not an optional
decoration on this loop: it determines which observable traces are legal.

\section{A state machine gives stopping a meaning}
An informal flag named \texttt{done} cannot explain whether a request succeeded,
was cancelled or failed. We need a state machine whose transitions are events.

\EngineFigure{request-states}{The request has one terminal owner. The loop can
emit zero or more ordered text events; its outgoing terminal transitions are
mutually exclusive. A failed request may already have emitted a prefix.}{fig:request-states}

Call the terminal outcome $o$. For a completed run with a working, non-panicking
in-process sink, our observable contract is
\begin{equation}
N_{\mathrm{terminal}}=1,\qquad
N_{\mathrm{text}}\geq 0,\qquad
\text{no output event follows the terminal event}.
\label{eq:terminal-contract}
\end{equation}
This is a contract for the teaching runtime's event stream, not a promise of
exactly-once network delivery. A disconnected client, process crash or retried
HTTP request introduces different failure domains.

\begin{proposition}[A single terminal owner simplifies the proof]
If every loop exit returns an outcome to one non-reentrant finalization path,
and that path emits one terminal event and then returns, the runtime cannot
emit two terminal events during an ordinary completed invocation.
\end{proposition}
\begin{proof}
Each loop exit transfers control out of the loop. Only the finalization path
contains the terminal emission. It executes once and returns without re-entering
the loop. Therefore there is one terminal emission, regardless of which exit
condition selected the outcome. This argument assumes ordinary control flow;
it does not claim recovery from process termination or a panicking callback.
\end{proof}

Why prefer this structure? If the sampler, model adapter and stream writer
each send their own final event, every pair of error paths must be checked for
double termination. Returning outcomes to one owner reduces the number of
places in which the invariant can be violated.

\section{Change one boundary at a time}
Keep the prompt fixed and run these experiments:
\begin{verbatim}
cargo run -q -p engine0 -- --max-tokens 1 'I like'
cargo run -q -p engine0 -- --cancel-at 1 'I like'
cargo run -q -p engine0 -- --fail-at 1 'I like'
\end{verbatim}
All three can expose the same generated prefix, a space followed by
\texttt{Rust}, but they have
different meanings. The first satisfies a requested output budget. The second
observes cancellation after one generated text token. The third injects a model
failure on the next invocation. An application that stores only the text cannot
reconstruct which experiment happened.

\begin{center}
\begin{tabular}{lll}
\toprule
Experiment & Visible prefix & Terminal class\\
\midrule
Default greedy & \texttt{Rust} & End of sequence\\
Budget one & \texttt{Rust} & Maximum tokens\\
Cancel at step one & \texttt{Rust} & Cancelled\\
Fail at step one & \texttt{Rust} & Failed\\
\bottomrule
\end{tabular}
\end{center}

The order of checks is observable. If both cancellation and the budget are
satisfied at the start of the next iteration, this runtime checks cancellation
first. Another engine could choose a different documented precedence. Neither
choice should be left to an accidental race between independent finalizers.

The implementation lives in \texttt{crates/engine0/src/lib.rs}; its
\texttt{Runtime::generate} method owns the loop. The lifecycle tests are in
\texttt{crates/engine0/tests/lifecycle.rs}. The textbook gate checks the commands
above against the executable, including their terminal classes and event counts.
The next chapters will reveal the tokenizer, model and sampler behind those
interfaces instead of changing the observable contract without warning.

\section{Follow three objects, not one vague arrow}
The request gives us three different journeys.

First, follow the \emph{token}. The prompt becomes IDs $[1,2]$, a model invocation
produces four scores, selection chooses ID 3, and ID 3 joins the next history.
The EOS ID is interpreted as control. A token is a vocabulary identity even
when its spelling happens to look like an English word.

Second, follow the \emph{byte}. Parameters occupy CPU memory. A forward call
reads them to compute temporary scores. On output, the tokenizer maps a
selected identity to a byte piece. A UTF-8 decoder may have to hold incomplete
bytes before emitting valid text. A token boundary is therefore not necessarily
a text-event boundary. \Appref{tokens} makes that failure concrete.

Third, follow the \emph{owner}. Immutable parameters can be shared conceptually;
generation history, decoder buffer, random state and terminal state belong to
one request. Sharing the weights does not grant permission to share mutable
history. The current engine executes synchronously, but choosing honest
ownership now gives later concurrency a tractable contract.

\EngineFigure{parameter-owners}{Two hypothetical requests read the same immutable
parameters while retaining distinct histories and activations. Dashed edges
are read-only borrows, not copies or implemented concurrent execution.}{fig:parameter-owners}

Training changes model parameters according to a learning rule. Inference
normally reads the fixed parameters while changing request-local state.
\emph{The weights do not change} does not mean \emph{nothing is mutable}. The
history grows, the decoder buffers bytes, counters advance and a sampled
request advances its own pseudorandom state.

\section{Measure a distance between named events}
Suppose an instrumented request records arrival $t_a$, execution start $t_s$,
first emitted token $t_k$, first text $t_b$ and terminal event $t_f$. These
timestamps must share a clock domain before we subtract them. Define
\begin{align}
T_{\mathrm{wait}}&=t_s-t_a,\\
T_{\mathrm{first\ token}}&=t_k-t_a,\\
T_{\mathrm{first\ text}}&=t_b-t_a,\\
T_{\mathrm{request}}&=t_f-t_a.
\end{align}
Waiting, computation and byte assembly now have different observable
boundaries. A request that selects EOS immediately has no first emitted text
token; its first-token and first-text measurements are absent, not zero.

\EngineFigure{request-timeline}{An illustrative clock, not a benchmark.
The first-token and first-text distances begin at arrival but end at different
events. A tokenizer can delay text until a byte sequence is complete.}{fig:request-timeline}

For the synthetic timestamps $(t_a,t_s,t_k,t_b,t_f)=(0,3,11,14,25)$ milliseconds,
the four distances are $3$, $11$, $14$ and $25$ milliseconds. The three
milliseconds between first token and first text do not prove that model
execution slowed down; they may be time spent waiting for another byte piece.

Likewise, aggregate output tokens per second are not the inverse of one
request's first-token latency. A batched service may process many requests
together, improving aggregate work per unit time while increasing an
individual request's wait. We will not call a measured throughput increase a
latency improvement without measuring that latency.

\section{Where the small engine stops being enough}
Our runtime rebuilds a history and calls a tiny CPU model synchronously.
There is no scheduler, persistent KV state, socket, retry protocol or device
transfer. Those omissions establish later problems rather than hiding them.

When many requests arrive, admission must decide which can fit. When prompts
have different lengths, a scheduler must balance prompt work with incremental
decode work. When a model uses attention, history-related state consumes
memory; reusing it requires layout, lifetime and validity rules. When output
crosses a network, a slow receiver creates backpressure and a disconnect
requires cancellation to propagate through the running computation.

The terminology should grow with those mechanisms. A \emph{model artifact}
is persistent data. A \emph{loaded model} is that data interpreted in a runtime.
An \emph{engine} executes and controls requests. A \emph{server} exposes an
interface. An operated \emph{service} adds deployment, resource limits,
observability and failure policy. Packaging the first as a file does not
implement the other four.

Production comparisons in later chapters are source-pinned case studies,
not substitutes for this construction. A repository containing a paged backend
does not prove the default request path uses it. A supported tensor type does
not prove a whole model family works. We will follow actual calls and data
when making those claims.

\section{What we have established}
We can now explain why one request needs both numerical prediction and a
lifecycle. We have observed a real selected token, feedback into another
prediction, an immediate-EOS counterexample and three distinct ways to stop
after the same visible prefix. We have named ownership and timing boundaries.

How text becomes an ID, and how an ID can fail to become readable text
immediately, is the subject of \Appref{tokens}, which opens the tokenizer and
the UTF-8 buffer without changing the request's terminal contract.

\LedgerSection

Streaming output as events and ending every request with exactly one terminal
outcome:

\begin{ConstraintLedger}
Latency & buys & The first words after the time to first token (0.74~s on the M1) instead of after the whole answer (6~s). \\
Memory capacity & spends & A batch slot and the request's KV cache from admission to its terminal event; an uncancelled orphan holds them to completion. \\
Compute & risks & Tokens generated after a client leaves are wasted: 490 in the opening's orphan (equation~\eqref{eq:orphan-waste}). \\
Correctness & buys & One terminal outcome per request, so a prefix is never mistaken for a completed answer. \\
\end{ConstraintLedger}

\LabSection{the request trace}

\LabIntent{In \texttt{code/labs/ch14-request-lifecycle/}, the mini-engine's
request trace (\texttt{engine0 --trace}) is extended with a streaming sink that
records every event with a monotonic timestamp and a cancellation source that
can fire at any declared boundary. Property tests drive thousands of random
event orders and check the terminal contract of
equation~\eqref{eq:terminal-contract}; the reader then adds a slow consumer and
measures how long a cancellation takes to reach the model loop.}

\HappenedSection{cancellation that did not reach the model}

The orphaned request of this chapter's opening was not an accident of the
measurement. llama.cpp's server, at the commit of the measured build, handles a
non-streamed request on an HTTP thread that waits for the request's result. It
waits in slices of one second, and only when a slice expires with no result
does it ask the socket whether the client is still there; if the client has
gone, it posts a cancellation that releases the request's slot
(\texttt{HTTP\_POLLING\_SECONDS} in \texttt{tools/server/server-context.cpp}
and \texttt{recv\_with\_timeout} in \texttt{tools/server/server-queue.cpp} at
commit \texttt{d006858}). The waiting threads of all requests share one
condition variable, and every result delivered for \emph{any} request wakes
them all; a woken thread that finds nothing for itself starts a fresh
one-second wait. While another client is streaming --- a result every 40~ms ---
no thread's wait ever expires, so no thread ever checks its socket. The
measurement showed exactly that: cancelled within a second and a half on a
quiet server, never cancelled while a neighbor streamed.

The first attempt to measure this found it by accident. The measuring script
polled the server's slot-status endpoint every 20~milliseconds to see when the
request would be released; each poll produced a result, each result reset the
waiting thread's timer, and the observer kept alive the request it was
observing. The measurement was rerun with a poll every 1.5 seconds, longer
than the server's slice, and with a separate streaming client as the
neighbor --- which is when the mechanism became clear. Both runs are in the
record.

vLLM makes the opposite design choice. Its OpenAI-compatible handlers are
wrapped so that each request's work races a listener for the HTTP server's
disconnect message; whichever finishes first cancels the other, and the
engine aborts the request when its generator is cancelled
(\texttt{vllm/entrypoints/serve/utils/api\_utils.py} and
\texttt{vllm/v1/engine/async\_llm.py} at commit \texttt{bcdacfc}). Disconnects
there are events, not polls. The lesson is the chapter's: cancellation is part
of the request's contract, and a contract enforced by a timer that other
traffic can reset is not enforced.

\FrontierSection{2026-09-24}

The lifecycle contract is becoming explicit in APIs. OpenAI's Responses API
streams typed events --- a response created, text deltas, a completion or an
error --- rather than bare tokens \cite{openai2026streaming}, and its
background mode decouples a request from the connection that started it:
a long response keeps running when the client disconnects, is cancelled only
by an explicit, idempotent cancel call, and can be resumed from a sequence
number if the stream drops \cite{openai2026background}. Engines are moving
the same way: detokenization and streaming in processes of their own
(\Chref{continuous-batching}) and cancellation driven by events.

\section{Worked problems: a request is more than its text}
\begin{minipage}{\linewidth}
\subsection{One prefix, four outcomes}
\textbf{Problem.} A log contains only the generated bytes for a space followed
by \texttt{Rust}. Can you infer that the model selected EOS? What minimum
additional evidence distinguishes the four experiments?

\textbf{Solution.} No. EOS, the token budget, cancellation and a failure on
the next invocation can all follow that prefix. Record the request identity,
ordered event types and final outcome. To distinguish a selected EOS from a
budget stop, record the selection trace as well: the budget may prevent the
next model invocation entirely. Text alone is not a completion certificate.
\end{minipage}\par

\begin{minipage}{\linewidth}
\subsection{Absent is not zero}
\textbf{Problem.} The prompt \texttt{I} produces an immediate EOS and no text.
What should its first-text latency field contain?

\textbf{Solution.} No value. The event defining the endpoint never happened.
Zero would assert that text was emitted at arrival. Including zeros in an
average would make a workload with many empty completions appear faster.
Report the count of defined measurements and the count of empty completions.
\end{minipage}\par

\Needspace{12\baselineskip}
\begin{minipage}{\linewidth}
\subsection{Test a cancellation boundary}
\textbf{Problem.} Design a test for cancellation observed after a model call
returns but before its token is emitted.

\textbf{Solution rubric.} Use a controlled cancellation source and a model
whose invocation is observable. Assert one invocation, zero subsequent token
or text events, exactly one Cancelled terminal event and no later events.
The test must not rely on sleeping for a guessed duration. Such a test proves
a cooperative boundary, not interruption of a running kernel.
\end{minipage}\par

\Needspace{8\baselineskip}
\begin{minipage}{\linewidth}
\subsection{Price an abandoned cache}
\textbf{Problem.} Illustratively, each abandoned request holds a constant
576~MiB for 37~s. Abandonments arrive steadily at 0.5 per second.
Find the occupancy per request and mean resident memory. What fails if the
cache grows during decode?

\textbf{Solution.} The memory-time cost is
$576\times37=21{,}312$~MiB\,s, or 20.8125~GiB\,s per request.
There are $0.5\times37=18.5$ abandoned requests resident on average, costing
$18.5\times576/1024=10.40625$~GiB. With growing caches, replace the constant
product by the integral of resident bytes over the abandoned lifetime.
Do not multiply by a request fraction and call it a fleet utilization
fraction without matching duration and resource populations.
\end{minipage}\par

\Needspace{8\baselineskip}
\begin{minipage}{\linewidth}
\subsection{Why a full queue cannot carry its own rescue}
\textbf{Problem.} Data enters at 25 events/s and leaves at five, with
capacity 64. Estimate its fill time and payload bytes at 306 bytes/event.
Should cancellation wait to enqueue behind that data?

\textbf{Solution.} Net growth is 20 events/s, giving 3.2~s from empty.
The payload occupies $64\times306=19{,}584$ bytes, excluding representation
overheads. A cancellation sent through a full data queue can never arrive
if the consumer has gone. The teaching owner instead has an independent
stop operation and stores terminal metadata outside the data deque.
Real services should also bound serialized bytes, not only event count.
\end{minipage}\par

\Needspace{8\baselineskip}
\begin{minipage}{\linewidth}
\subsection{A terminal before retirement}
\textbf{Problem.} The owner submits ticket $(7,12,0)$, commits Cancelled
and delivers its terminal. The worker still holds the ticket. May slot seven
be reassigned? What should happen to its late reply?

\textbf{Solution.} Not yet. The terminal fixes output policy, not device
completion. A matching retirement clears the outstanding ticket and drops
the late payload. Only then is the modeled lease reusable. If generation
13 later receives a duplicate reply from generation 12, it must reject it
without clearing generation 13's ticket. A same-generation stale serial
must also fail. The lab tests both cases; a production device fence and
shared-block reference accounting remain separate obligations.
\end{minipage}\par

\Needspace{8\baselineskip}
\begin{minipage}{\linewidth}
\subsection{A wake is not our result}
\textbf{Problem.} An illustrative handler starts at 0~ms, waits in 100~ms
slices, and restarts after unrelated wakes every 4~ms through 2,000~ms.
When is its first timeout? What changes if its deadline was fixed at 100~ms?

\textbf{Solution.} The final wake establishes expiry at
$2{,}000+100=2{,}100$~ms. A client gone at 350~ms is therefore unobserved
for about 1,750~ms. A fixed initial deadline stays at 100~ms: at the 4~ms
wake, 96~ms remain. A periodic detector needs subsequent fixed checks to
notice the later disconnect; the first 100~ms check precedes it.
These are chosen integer times, not an operating-system timing guarantee.
\end{minipage}\par

\Needspace{8\baselineskip}
\begin{minipage}{\linewidth}
\subsection{Falsify a weak cancellation checker}
\textbf{Problem.} A three-token-budget request produces tokens 0, 1, 2 and
one MaxTokens terminal. The harness records a cancellation while the
request was still running after token 0. Why is checking the natural
outcome insufficient?

\textbf{Solution.} The counts match the uncancelled oracle but violate the
observed cancellation contract. In the step engine, cancellation precedes
the next numerical step, so the legal terminal is Cancelled after one token.
A checker must reject the natural trace, not accept it because it ends
once. This is the negative test
\texttt{checker\_rejects\_natural\_end\_after\_observed\_cancellation}.
Safety, natural stopping and cooperative cancellation are different checks.
\end{minipage}\par

